% BHU PYQ -- Auto-generated & error-corrected
\documentclass[11pt,a4paper]{article}

\usepackage[a4paper,top=2.5cm,bottom=2.5cm,left=2.8cm,right=2.8cm,headheight=14pt]{geometry}
\usepackage{amsmath,amssymb}
\usepackage[T1]{fontenc}
\usepackage[utf8]{inputenc}
\usepackage{lmodern}
\usepackage{microtype}
\usepackage{fancyhdr}
\usepackage{enumitem}
\usepackage{hyperref}
\usepackage{lastpage}
\usepackage{parskip}

\hypersetup{colorlinks=true,urlcolor=black,linkcolor=black,
  pdftitle={Paper-IX: Statistical Inference and Decision Theory (Old Course) -- BHU PYQ},pdfauthor={Banaras Hindu University}}

\pagestyle{fancy}\fancyhf{}
\renewcommand{\headrulewidth}{0.4pt}\renewcommand{\footrulewidth}{0.4pt}
\fancyhead[L]{\small\textit{Banaras Hindu University}}
\fancyhead[R]{\small\textit{Paper-IX: Statistical Inference \& Decision Theory}}
\fancyfoot[C]{\small Page \thepage\ of \pageref{LastPage}}

\newlist{parts}{enumerate}{2}
\setlist[parts,1]{label=(\alph*),leftmargin=*,itemsep=5pt,topsep=3pt}
\setlist[parts,2]{label=(\roman*),leftmargin=*,itemsep=3pt,topsep=2pt}
\newcommand{\pts}[1]{\hfill\textit{[#1]}}

\begin{document}

\begin{center}
  {\large\bfseries BANARAS HINDU UNIVERSITY}\\[2pt]
  {\normalsize B.Sc. (Hons.) Part-III Examination, 2014-15}\\[6pt]
  \rule{0.8\linewidth}{0.5pt}\\[6pt]
  {\large\bfseries Paper No. Paper-IX: Statistical Inference \& Decision Theory (Old Course)}\\[4pt]
  {\normalsize Time: 3 Hours\hfill Maximum Marks: 70}
\end{center}
\rule{\linewidth}{0.4pt}
\smallskip
\noindent\textit{Instructions: Question No.\ 1 is compulsory. Attempt any \textbf{four} questions, selecting one from each Unit. Symbols have their usual meanings.}
\bigskip

\medskip\noindent\textbf{Question 1 (Compulsory).} Answer all parts briefly:\pts{20}

\begin{parts}
  \item What is the difference between an estimator and an estimate? \pts{2}
  \item State the invariance property of a consistent estimator. \pts{2}
  \item Define a complete statistic. \pts{2}
  \item When is a statistical test said to be unbiased? \pts{2}
  \item Define level of significance. \pts{2}
  \item Define the Minimum Variance Bound (MVB) estimator. \pts{2}
  \item Define confidence limits and confidence coefficient. \pts{2}
  \item Define the minimax decision rule. \pts{2}
  \item Define a statistical parameter with an example. \pts{2}
  \item Define the power function of a statistical test. \pts{2}
\end{parts}

\medskip\rule{\linewidth}{0.3pt}
\begin{center}\textbf{UNIT I}\end{center}

\medskip\noindent\textbf{Question 2.}
\begin{parts}
  \item Define consistency of an estimator. Show that an estimator $T_n$ of parameter $\theta$ based on a sample of size $n$ is consistent if:
  \[
  E(T_n) \to \theta \quad \text{and} \quad \mathrm{Var}(T_n) \to 0 \quad \text{as } n \to \infty.
  \]\pts{7}
  \item What are the desirable properties of a good estimator? Explain any one of them in detail with a suitable example. \pts{7}
\end{parts}

\medskip\rule{\linewidth}{0.3pt}
\begin{center}\textbf{UNIT II}\end{center}

\medskip\noindent\textbf{Question 3.}
\begin{parts}
  \item Define a sufficient statistic and state the Fisher--Neyman factorization theorem. \pts{7}
  \item State and prove the Rao--Blackwell theorem. \pts{7}
\end{parts}

\medskip\rule{\linewidth}{0.3pt}
\begin{center}\textbf{UNIT III}\end{center}

\medskip\noindent\textbf{Question 4.}
\begin{parts}
  \item Explain the method of interval estimation. Obtain the $95\%$ confidence interval for the variance $\sigma^2$ of a normal distribution $N(\mu, \sigma^2)$ when $\mu$ is unknown. \pts{7}
  \item Explain the method of maximum likelihood estimation and state the asymptotic properties of maximum likelihood estimators (MLEs). \pts{7}
\end{parts}

\medskip\noindent\textbf{Question 5.}
\begin{parts}
  \item Explain the method of moments for estimating parameters. Discuss its merits and demerits compared to the method of maximum likelihood. \pts{7}
  \item Define and explain simple and composite hypotheses with appropriate examples. \pts{7}
\end{parts}

\medskip\rule{\linewidth}{0.3pt}
\begin{center}\textbf{UNIT IV}\end{center}

\medskip\noindent\textbf{Question 6.}
\begin{parts}
  \item Explain the following terms in statistical decision theory:
  \begin{parts}
    \item Loss function
    \item Risk function
    \item Bayes' decision rule
  \end{parts}\pts{7}
  \item State and prove the Neyman--Pearson lemma. Use it to obtain the most powerful test for testing $H_0\colon \mu = \mu_0$ against $H_1\colon \mu = \mu_1\;(\mu_1 > \mu_0)$ on the basis of a sample of size $n$ from $N(\mu, \sigma^2)$ when $\sigma^2$ is known. \pts{7}
\end{parts}

\medskip\noindent\textbf{Question 7.}

Describe the procedure of the Likelihood Ratio Test (LRT). Obtain the likelihood ratio test for testing $H_0\colon \mu = \mu_0$ against $H_1\colon \mu \ne \mu_0$ on the basis of a sample of size $n$ from $N(\mu, \sigma^2)$ when $\sigma^2$ is unknown. \pts{14}

\vfill
\rule{\linewidth}{0.4pt}
\begin{center}\small
  \href{https://bhu-exam.vercel.app/}{Click here to visit our website}\\[4pt]
  \href{https://me-aryan.vercel.app/}{Click here to visit our contributor}\\[4pt]
  \href{https://quashan-me.vercel.app/}{Click here to visit our contributor}
\end{center}

\end{document}